Large language models are increasingly used for ideation, recommendation, and creative tasks---but can they identify what is genuinely underrated, or do they simply echo the crowd's contrarianism? We introduce the ``most underrated'' prompt as a diagnostic for generative novelty. Querying three OpenAI models (\gptfull, \gptmini, \gptfo) with 80 underrated prompts across 8 categories (4,800 responses total), we find strong intra-model convergence: the most popular answer captures 48.7\% of responses on average, with extreme cases like ``Children of Men'' (most underrated movie) reaching 87--95\% dominance across all models. These top answers---sumac, the Sega Dreamcast, Lise Meitner---are well-known ``underrated'' picks that recur in online discussions, revealing a \metapop where models learn what is \emph{popular as underrated} rather than what is genuinely overlooked. Surprisingly, inter-model agreement is low (8.8\% top-1 unanimity), indicating that different models have learned \emph{different} but equally stereotyped notions of underratedness. Compared to ``best'' and factual control prompts, underrated prompts elicit more lexical diversity (unique rate 0.40 vs.\ 0.21 vs.\ 0.13) but this reflects a wider answer space, not genuine novelty. Our results demonstrate that current LLMs possess a narrow, training-data-derived sense of ``taste'' rather than true discernment, and we propose the underrated prompt as a lightweight, reproducible diagnostic for evaluating generative novelty in future models.
